\documentclass[11pt,letterpaper]{article}
\usepackage[margin=1in]{geometry}
\usepackage{amsmath,amssymb,amsthm}
\usepackage{hyperref}
\usepackage{cleveref}
\usepackage{booktabs}

\newtheorem{theorem}{Theorem}[section]
\newtheorem{lemma}[theorem]{Lemma}
\newtheorem{proposition}[theorem]{Proposition}
\newtheorem{corollary}[theorem]{Corollary}
\newtheorem{definition}[theorem]{Definition}

\numberwithin{equation}{section}
\renewcommand{\theequation}{S13.\arabic{equation}}

\title{Supplement S13: CKM Matrix from Vertex Partition \\
\large Supporting Manuscript \S16}
\author{UDT Collaboration}
\date{}

\begin{document}
\maketitle

\begin{abstract}
We derive the four CKM parameters from the angular quantum numbers
$(j, \ell, |\kappa_\mathrm{max}|) = (1/2, 1, 3)$.  The Cabibbo angle
$\sin\theta_C = 9/40$ ($-0.13\%$) arises from the vertex partition,
the Wolfenstein parameter $A = \cos(\pi/5)$ ($-0.24\%$) from the
metric depth, and the CP-violating phase
$\delta = \arctan(9/4)$ ($0.1\sigma$) encodes the force hierarchy
$\alpha_s/\alpha_\mathrm{EM} = 9/4$.  The strong CP problem is
solved: $\theta_\mathrm{QCD} = 0$ from the reality of the $S^2$
angular algebra.  Verified by \texttt{analysis/19\_ckm\_mixing.py}.
\end{abstract}

\tableofcontents

%----------------------------------------------------------------------
\section{Bilinear vs.\ vertex partition}
\label{sec:partition}
%----------------------------------------------------------------------

The PMNS angles (Supplement~S7) arise from \emph{bilinear} partitions:
two-field operators $\bar{\psi}\psi$ with squared multiplicities and
denominator $13 = (2j+1)^2 + (2\ell+1)^2$.  The CKM angles arise from
\emph{vertex} partitions: three-field operators $\bar{\psi}A\psi$ with
cubed spin and denominator $40 = (2j+1)^3(2|\kappa_\mathrm{max}|-1)$.

This distinction explains the PMNS--CKM hierarchy:
\begin{itemize}
  \item PMNS: $\sin^2\theta_{12} = 4/13 \approx 0.31$ (large).
  \item CKM: $\sin\theta_C = 9/40 = 0.225$ (small).
\end{itemize}
Bilinear partitions are $O(1)$ fractions; vertex partitions are
suppressed by the additional $(2j+1)$ factor in the denominator.

%----------------------------------------------------------------------
\section{The Cabibbo angle}
\label{sec:cabibbo}
%----------------------------------------------------------------------

At a three-field vertex $\bar{\psi}A\psi$, each field carries a
$(2j+1)$ multiplicity.  The vertex partition is the fraction occupied
by the orbital sector:
\begin{equation}\label{eq:cabibbo}
  \sin\theta_C = \frac{(2\ell+1)^2}{(2j+1)^3(2|\kappa_\mathrm{max}|-1)}
  = \frac{9}{40} = 0.2250\,,
\end{equation}
compared to the experimental value $0.2253$ ($-0.13\%$).

The numerator $(2\ell+1)^2 = 9$ counts the orbital coupling channels
(the same factor appearing in the polarization part of $1/\alpha$).
The denominator $(2j+1)^3 = 8$ cubes the spin multiplicity (three
fields at the vertex), and $(2|\kappa_\mathrm{max}|-1) = 5$ is the
angular content.

%----------------------------------------------------------------------
\section{The Wolfenstein parametrization}
\label{sec:wolfenstein}
%----------------------------------------------------------------------

The Wolfenstein parameters are:
\begin{align}
  \lambda &= \sin\theta_C = \frac{9}{40} = 0.2250
    \qquad(-0.13\%)\,, \\
  A &= |\phi_0| = \cos(\pi/5) = 0.8090
    \qquad(-0.24\%)\,.
\end{align}
The parameter $A$ is the metric depth itself --- the same geometric
constant that determines the cavity well depth and the Airy finesse
$F = 16/9$ in the CMB (Manuscript \S21).

\subsection{CKM matrix elements}

The first two rows of the CKM matrix, to $O(\lambda^4)$:
\begin{equation}
  V = \begin{pmatrix}
    1 - \lambda^2/2 & \lambda & A\lambda^3 \\
    -\lambda & 1 - \lambda^2/2 & A\lambda^2 \\
    \cdots & \cdots & \cdots
  \end{pmatrix}\,,
\end{equation}
with $\lambda = 9/40$ and $A = \cos(\pi/5)$.  Seven elements are
determined at sub-percent accuracy from these two quantities.

The remaining Wolfenstein parameters follow from
$\bar{\rho} = 1/(2\pi)$ (angular normalization per radian, $0.02\sigma$)
and $\bar{\eta} = 9/(8\pi)$ ($1.0\sigma$), with
$\bar{\eta}/\bar{\rho} = 9/4 = \tan\delta$ automatic from the CP phase.
The magnitude $R_b = \sqrt{\bar{\rho}^2 + \bar{\eta}^2}
= \sqrt{97}/(8\pi) = 0.392$.

%----------------------------------------------------------------------
\section{CP violation as the force hierarchy}
\label{sec:cp}
%----------------------------------------------------------------------

The CKM CP-violating phase is:
\begin{equation}\label{eq:delta-ckm}
  \delta_\mathrm{CKM} = \arctan\!\frac{9}{4}
  = \arctan\!\frac{(2\ell+1)^2}{(2j+1)^2}
  = 66.04^\circ\,,
\end{equation}
compared to the experimental value $65.5^\circ \pm 2.5^\circ$
($0.1\sigma$).

The argument $9/4 = \alpha_s/\alpha_\mathrm{EM}$
(\cref{eq:alpha-ratio} of the main text): the tangent of the
CP-violating phase \emph{is} the ratio of the strong to
electromagnetic coupling constants.  CP violation in the quark sector
is not an independent phenomenon --- it is the force hierarchy
expressed as a phase.

\subsection{Jarlskog invariant}

The rephasing-invariant measure of CP violation:
\begin{equation}
  J_\mathrm{CKM} = \mathrm{Im}(V_{us}V_{cb}V_{ub}^*V_{cs}^*)
  \approx A^2\lambda^6\sin\delta
  = 3.03 \times 10^{-5}\,,
\end{equation}
compared to the experimental value $3.00 \times 10^{-5}$ ($+0.8\%$).

%----------------------------------------------------------------------
\section{The strong CP problem}
\label{sec:strong-cp}
%----------------------------------------------------------------------

The QCD Lagrangian admits a CP-violating $\theta$-term:
$\mathcal{L}_\theta = \theta\,g_s^2/(32\pi^2)\,G\tilde{G}$.
Experiment constrains $|\theta_\mathrm{QCD}| < 10^{-10}$.  The
Standard Model offers no explanation for this smallness.

In UDT, the angular algebra on $S^2$ that closes as
$\mathfrak{su}(3)$ (Supplement~S16) has \emph{real} structure
constants.  The tensor operators $T_a$ are Hermitian with real
commutation relations.  A $\theta$-term would require an imaginary
component in the structure constants, which the $S^2$ algebra does
not admit.  Therefore:
\begin{equation}
  \theta_\mathrm{QCD} = 0 \qquad\text{(exact)}\,.
\end{equation}
No axion, no Peccei--Quinn symmetry, no fine-tuning.  The strong CP
problem is solved by the reality of the angular geometry.

%----------------------------------------------------------------------
\section{Unified mixing picture}
\label{sec:unified}
%----------------------------------------------------------------------

\begin{center}
\begin{tabular}{lccl}
\toprule
Sector & Type & Denominator & Angles \\
\midrule
PMNS & Bilinear ($\bar\psi\psi$) & $13, 7, 45$
  & Large ($O(0.3\text{--}0.6)$) \\
CKM & Vertex ($\bar\psi A\psi$) & $40$
  & Small ($O(0.2)$) \\
\bottomrule
\end{tabular}
\end{center}

Six fundamental mixing parameters + two CP phases from three quantum
numbers.  The bilinear/vertex distinction is geometric: it counts
how many fields participate in the coupling.

%----------------------------------------------------------------------
\section{Verification}
\label{sec:verification}
%----------------------------------------------------------------------

All results verified by \texttt{analysis/19\_ckm\_mixing.py}.
Output: \texttt{data/generated/19\_ckm\_mixing.json}.

\end{document}
